\subsection*{12.4 Application. MMSE Estimation}

The minimum mean-squared error (MMSE) criterion is a widely used approach for

optimal linear filtering in signal processing. The Orthogonality Principle provides a

powerful tool for deriving the linear MMSE estimator. For more details, readers can

refer to textbooks on estimation theory, such as [ 5]In this section, we will present

some examples of linear and nonlinear MMSE estimation.

12.4.1 Linear MMSE Estimator

Suppose Y , X1,X2,...,X n are inL2(P) We want to find a1,a 2,...,a n such that

the mean-squared error

Y-

n\sum

i=1

aiXi2

is minimized.

Let W be the space consisting of all linear combinations .

\sumn

i=1 aiXi One can

check that this is a closed subspace. By the Orthogonality Principle in Theorem12.5,

we can solve for ai from

\langleY-

n\sum

j= 1

aj Xj ,Xi\rangle= 0,for i = 1 ,2,...,n .

This amounts to solving a system of linear equations:

E[X1X1]E [X1X2]\cdot\cdot\cdot E[X1Xn]

E[X2X1]E [X2X2]\cdot\cdot\cdot E[X2Xn]

..

.. ... ...

E[XnX1]E [XnX2]\cdot\cdot\cdot E[XnXn]

a1

a2

..

an

=

E[X1Y]

E[X2Y]

..

E[XnY]

When the matrix on the left is nonsingular, we can solve for ai 's. We remark that

whenX1,...,X n have zero mean, this matrix is the same as the covariance matrix.

\textbf{Example 12.4.1 (Linear MMSE Estimator with Two Input Random Variables)} \\

When n= 2 , the problem is to find a1 and a2 that minimize Y- a 1X1 -a 2X22. The optimal

solution is given by the following system of linear equations:

[

E[X2

1]E [X1X2]

E[X1X2]E [X2

2]

][

a1

a2

]

=

[

E[X1Y]

E[X2Y]

]

The matrix on the left is positive semi-definite by the Cauchy-Schwarz inequality. The determinant

is positive if and only if X1 and. X2 are linearly independent.

We can use this system of linear equations to solve the linear regression problem in Example12.2.3,

by substituting the expectation E[X2

1] by . 1

n

\sumn

i=1 x2

i , the inner product E[X1X2] by . 1

n

\sumn

i=1 xi ,

etc. By solving

[ 1

n

\sumn

i=1 x2

i

n

\sumn

i=1 xi

n

\sumn

i=1 xi 1

][

a

b

]

=

[ 1

n

\sumn

i=1 xiyi

n

\sumn

i=1 yi

]

,

we recover the answer to the least square problem:

a= n \sum

i xiyi - \sum

i xi

\sum

i yi

n \sum

i x2 -( \sum

i xi)2 ,

b=

\sum

i yi -a \sum

i xi

n .

\textbf{Example 12.4.2 (Scaling Factor for Linear MMSE Estimate)} \\

Suppose we want to estimate a random variable X , which is observed with some additive noise N .

The observed random variable is Y= cX + N ,w h e r e c is a nonzero constant. We suppose that X

has mean \mu and variance \sigma2

X , N has zero mean and variance \sigma2

N ,a n d X and N are independent.

We also suppose that the parameters c ,\mu,\sigma2

X ,a n d. \sigma2

N are known, andc/=0 .

An unbiased estimate of X is simplyX= Y/c , if we can measure the random variable Y H o w e v e r ,

if we want to minimize the mean-squared error, the optimal constant k that minimizes the mean-

squared errorkY- X 2 can be computed by differentiating

k2E[Y 2]- 2kE[XY]+ E[X2]

with respect to k The optimal value scaling coefficient is

k* = E[XY]

E[Y 2] = c\sigma2

X

c2\sigma2

X +\sigma 2

N

The linear MMSE estimate of X is then given by

XMMSE = c2\sigma2

X

c2\sigma2

X +\sigma 2

N

\cdot Y

c .

We multiply the unbiased estimateY/c by the scaling factor.

c2\sigma2

X

c2\sigma2

X+\sigma2

N

to obtain the MMSE estimate,

which is a linear function of the observed variable Y .

12.4.2 Nonlinear MMSE Estimation

In general, given Y , X1, X2,...,X n \inL 2(P) , the problem of nonlinear MMSE

estimation is to find a measurable function g(x1,x 2,...,x n) such that

Y- g(X 1,X 2,...,X n)2

is minimized. We formulate the problem in terms of \sigma-algebra.

Recall that \sigma(X1,X 2,...,X n) denotes the smallest \sigma-algebra of \mathcal{F} such that

X1,...,X n are measurable. A function g(x1,x 2,...,x n) mapping from. Rn to. R is

measurable with respect to \sigma(X1,X 2,...,X n) iff

g- 1(B)\in \sigma(X 1,X 2,...,X n)for all B \in \mathcal{B} (R).

The next theorem characterizes the random variables that can be written as a

measurable function of X1 to. Xn.

\begin{thmbox}{12.6 (Functions Measurable with Respect to\sigma( X))}
\end{thmbox}

X()= (X 1(), X2(), . . . , Xn())

Consider an \mathcal{F}-measurable vector function X()=

(X1(), X2(), . . . , Xn()) defined on a measurable space .(\Omega,\mathcal{F})A n

\mathcal{F}-measurable function Y() can be written as Y= g(X 1,X 2,...,X n) for

a .(\mathcal{B}(Rn),\mathcal{B}(R))-measurable function g if and only if Y is measurable with

respect to\sigma( X).

\textit{Proof} ( . ) Suppose that X is \mathcal{F}-measurable from \Omega to Rn, and g is

(\mathcal{B}(Rn),\mathcal{B}(R))-measurable function from Rn to R. Consider the composite

function Y= g X Given any Borel set B\in \mathcal{B} (R), the inverse image g- 1(B) is

\mathcal{B}(Rn)-measurable. Since the vector function X is.(\sigma(X),\mathcal{B}(Rn))-measurable, the

inverse imageY- 1(B)= X - 1(g- 1(B)) is in\sigma( X).

(. ) Suppose that Y is an \mathcal{F}-measurable function that is also \sigma( X)-measurable.

By definition, we have Y- 1(B)\in \sigma( X) for all Borel sets B in. R.

When Y is an indicator function .1A() for some A\in \mathcal{F} , the inverse image

Y- 1({1})= A must be in \sigma( X) by the assumption that Y is \sigma( X)-measurable.

Hence,A= X - 1(B) for someB\in \mathcal{B} (Rn)We can set

g(x1,x 2,...,x n)=

{

1i f (x1,x 2,...,x n)\in B,

0 otherwise .

Then g is Borel measurable, and Y() =g( X()) for all\in \Omega .

Next, suppose Y is a simple function in the form .

\summ

i=1 ai 1Ai ,f o r s o m eAi \in\mathcal{F}.

We can find Borel measurable function gi such that

Y() =

m\sum

i=1

aigi(X1(), . . . , Xn()).

The rest of the proof is to extend it to nonnegative measurable functions and real-

valued measurable functions. It is the standard machinery of Lebesgue integral and

the detail is omitted. \hfill $\square$

In nonlinear MMSE estimation, we define W to be the set of all random variables

on .(\Omega,\mathcal{F},P) that are measurable with respect to the \sigma-algebra \sigma( X) and have

finite second moment. We denote this set by L2(\Omega, \sigma(X), P), where P denotes the

restriction of the original probability measure to the sub- \sigma-algebra\sigma( X)Note that

L2(\Omega, \sigma(X), P) is a closed subspace of the space of all square-integrable random

variables on.(\Omega,\mathcal{F},P) .

The optimization problem can now be stated as

min

g(x1,...,xn)

Y- g(X 1,X 2,...,X n)2

with the minimum taken over all functions g: R n \to R such that. g(X1,...,X n)\in

L2(\Omega, \sigma(X), P)We project Y onto the closed subspace L2(\Omega, \sigma(X), P)B y

Theorem 12.4, a solution exists and any two solutions are equal almost surely. By

the Orthogonal Principle, the solution is characterized by

E[ YZ ]= E[YZ ]for all Z \in L 2(\Omega, \sigma(X), P). (12.5)

Indeed, the condition in (12.5) can be written in several equivalent forms,

\langleY- Y,Z \rangle= 0 for all Z \in L 2(\Omega, \sigma(X), P)

E [(Y- Y)Z ]= 0 for all Z \in L 2(\Omega, \sigma(X), P)

E [YZ ]= E[ YZ ]for all Z \in L 2(\Omega, \sigma(X), P)

E[Yh( X)]= E[ Yh( X)] for all \mathcal{B}(Rn )-measurable h(x) st. h(X) \in L2(P).

We illustrate this idea in the following example.

\textbf{Example 12.4.3 (An Example of Nonlinear MMSE Estimation)} \\

Consider a sample space .{a,b,c,d } and define a uniform distribution on the four sample points.

Define two random variables using the table

a b c d

X(. ) 1 1 2 2

Y(. ) 1 2 3 4

We want to find the nonlinear MMSE estimator of Y given X. Since the probability space is finite,

all functions defined on this probability space are measurable and have finite second moments. The

problem is equivalent to finding a function g(x) that minimizesY - g(X)2.

The\sigma-algebra generated by X contains four sets \emptyset,\Omega ,.{a,b },a n d.{c,d }The MMSE estimator of

Y given X is a function of X and is in the form

Y() =

{

\alphaif = a or = b

\betaif = c or = d.

The optimal values of \alpha and \beta can be determined using calculus. The objective function that we

are going to minimize is a multi-variable quadratic function,

4(1- \alpha) 2 + 1

4(2- \alpha) 2 + 1

4(3- \beta) 2 + 1

4(4- \beta) 2.

The solution can be derived by partially differentiating the objective function with respect to \alpha and

\beta and solving a system of linear equations.

In the followings, we can obtain the answer using the Orthogonal Principle. The MMSE estimator

Y satisfies the following property:

E[YZ ]= E[ YZ ] for all \sigma(X)-measurable functions Z.

We first considerZ = 1{a,b}Applying (12.5),we get

E[Y 1{a,b}]= 1

4 + 2

4 +0 + 0 ,

E[ Y1{a,b}]= \alphaE[1{a,b}]= \alpha

( 1

4 + 1

)

The second line follows from the fact that Y is constant on the event .{a,b }Equating the above

two values yields\alpha = 1.5.

We next consider the function Z = 1{c,d}We have

E[Y 1{c,d}]= 0 + 0 + 3

4 + 4

4,

E[ Y1{c,d}]= \betaE[1{c,d}]= \beta

( 1

4 + 1

)

This gives\beta= 3 .5.

The optimal function g(x) that minimizes mean-squared error Y- g(X) 2 is

g(x)=

{

1.5i f x= 1 ,

3.5i f x= 2 .

We implement this example using Python.

from random import choice

Omega = ['a','b','c','d'] # sample space

Y = {'a':1, 'b':2, 'c':3, 'd':4} # represent Y by a dictionary

omega= choice(Omega) # pick a sample uniformly

X=1i f (Y[omega] <= 2) else 2 # X equals 1 if Y=1, 2 if Y=2

Y_hat = 1.5 if X==1 else 3.5 # estimate random variable Y by X

print(f"Y={Y[omega]}, X={X}, Y_hat={Y_hat}")

In the program, we denote the estimate of Y by Y_hat , which is a function of X A sample output

of this program is:

Y=4, X=2, Y_hat=3.5

The conditional expectation operator in L2(P) generally provides the nonlin-

ear minimum mean-squared error estimation. Specifically, the function g(x) that

minimizes the mean-squared error E[(Y- g(X)) 2] is given by the conditional

expectationg(x)= E [Y\vertX= x ].
